% Part: incompleteness
% Chapter: theories-computability
% Section: extensions-of-q-not-decidable

\documentclass[../../../include/open-logic-section]{subfiles}

\begin{document}

\olfileid{inc}{tcp}{cqn}

\olsection{امتدادات $\Th{Q}$ المتسقة غير قابلة للقرار}

\begin{explain}
الاتساق، كما تقدم، أن تكون النظرية \emph{متسقة} بمعنى أنها لا تثبت
$!A$ و$\lnot !A$ معًا، أيًّا كانت الصيغة $!A$. فالتناقض يستلزم كل
شيء، ولذلك تصير النظرية غير المتسقة تافهة: تثبت فيها كل جملة.
ومن البيّن أن الاتساق-$\omega$ يستلزم الاتساق، دون العكس؛ فثمة
نظريات متسقة ليست متسقة-$\omega$. فإذا اقتصرنا في فرض
\olref[oqn]{thm:oconsis-q} على الاتساق المعتاد، حصلنا على مبرهنة
أقوى، وهو ما سنفعله.
\end{explain}

\begin{lem}
يمتنع وجود «علاقة شاملة قابلة للحساب». والمراد أنه لا توجد علاقة
ثنائية قابلة للحساب ${\OLId{latin-looped}{R}}({\OLId{latin-ordinary}{x}},{\OLId{latin-ordinary}{y}})$ تستوفي هذا الشرط: لكل علاقة أحادية
قابلة للحساب ${\OLId{latin-looped}{S}}({\OLId{latin-ordinary}{y}})$ عدد ${\OLId{latin-ordinary}{k}}$ يجعل، لكل ${\OLId{latin-ordinary}{y}}$، صدق ${\OLId{latin-looped}{S}}({\OLId{latin-ordinary}{y}})$ مكافئًا لصدق
${\OLId{latin-looped}{R}}({\OLId{latin-ordinary}{k}},{\OLId{latin-ordinary}{y}})$.
\end{lem}

\begin{proof}
لنفرض وجود علاقة شاملة قابلة للحساب ${\OLId{latin-looped}{R}}({\OLId{latin-ordinary}{x}},{\OLId{latin-ordinary}{y}})$، ولنعرّف ${\OLId{latin-looped}{S}}({\OLId{latin-ordinary}{y}})$ بأنها
$\lnot {\OLId{latin-looped}{R}}({\OLId{latin-ordinary}{y}},{\OLId{latin-ordinary}{y}})$. حينئذ تكون ${\OLId{latin-looped}{S}}({\OLId{latin-ordinary}{y}})$ قابلة للحساب، فيوجد عدد ${\OLId{latin-ordinary}{k}}$ تكون
عنده ${\OLId{latin-looped}{S}}({\OLId{latin-ordinary}{y}})$ مكافئة لـ${\OLId{latin-looped}{R}}({\OLId{latin-ordinary}{k}},{\OLId{latin-ordinary}{y}})$. فتكون ${\OLId{latin-looped}{S}}({\OLId{latin-ordinary}{k}})$ مكافئة لـ${\OLId{latin-looped}{R}}({\OLId{latin-ordinary}{k}},{\OLId{latin-ordinary}{k}})$
ولنقيضها $\lnot {\OLId{latin-looped}{R}}({\OLId{latin-ordinary}{k}},{\OLId{latin-ordinary}{k}})$ معًا، وهذا هو التناقض.
\end{proof}


\begin{thm}
كل نظرية $\Th{T}$ متسقة تشتمل على $\Th{Q}$ تكون
$\Th{T}$ غير قابلة للقرار.
\end{thm}


\begin{proof}
لنفرض أن $\Th{T}$ امتداد متسق قابل للقرار لـ$\Th{Q}$.
وسنبني بواسطة $\Th{T}$ علاقة شاملة قابلة للحساب، فيلزم المحال.

نعرّف ${\OLId{latin-looped}{R}}({\OLId{latin-ordinary}{x}},{\OLId{latin-ordinary}{y}})$ بالشرط الآتي، اللازم والكافي لصدقها:
\begin{quote}
${\OLId{latin-ordinary}{x}}$ شفرة لصيغة $!D({\OLId{latin-ordinary}{u}})$، و$\Th{T}$ تثبت $!D(\num {\OLId{latin-ordinary}{y}})$.
\end{quote}
فمن فرض قابلية $\Th{T}$ للقرار تكون ${\OLId{latin-looped}{R}}$ قابلة للحساب. وأما شمول
${\OLId{latin-looped}{R}}$ فبيانه أن كل علاقة قابلة للحساب ${\OLId{latin-looped}{S}}({\OLId{latin-ordinary}{y}})$ تقبل التمثيل في
$\Th{Q}$، ومن ثم في $\Th{T}$، بصيغة $!D_{\OLId{latin-looped}{S}}({\OLId{latin-ordinary}{u}})$. فلكل ${\OLId{latin-ordinary}{n}}$:
\begin{eqnarray*}
{\OLId{latin-looped}{S}}({\OLId{latin-ordinary}{n}}) & \lif & {\OLId{latin-looped}{T}} \vdash !D_{\OLId{latin-looped}{S}}(\num {\OLId{latin-ordinary}{n}}) \\
& \lif & {\OLId{latin-looped}{R}}(\Gn{!D_{\OLId{latin-looped}{S}}({\OLId{latin-ordinary}{u}})}, {\OLId{latin-ordinary}{n}})
\end{eqnarray*}
وكذلك
\begin{eqnarray*}
\lnot {\OLId{latin-looped}{S}}({\OLId{latin-ordinary}{n}}) & \lif & {\OLId{latin-looped}{T}} \vdash \lnot !D_{\OLId{latin-looped}{S}}(\num {\OLId{latin-ordinary}{n}}) \\
& \lif & {\OLId{latin-looped}{T}} \not\vdash !D_{\OLId{latin-looped}{S}}(\num {\OLId{latin-ordinary}{n}}) \quad \text{(لأن $\Th{T}$
  متسقة)} \\
& \lif & \lnot {\OLId{latin-looped}{R}}(\Gn{!D_{\OLId{latin-looped}{S}}({\OLId{latin-ordinary}{u}})}, {\OLId{latin-ordinary}{n}}).
\end{eqnarray*}
وبذلك يكون، لكل ${\OLId{latin-ordinary}{y}}$، صدق ${\OLId{latin-looped}{S}}({\OLId{latin-ordinary}{y}})$ مكافئًا لصدق
${\OLId{latin-looped}{R}}(\Gn{!D_{\OLId{latin-looped}{S}}({\OLId{latin-ordinary}{u}})},{\OLId{latin-ordinary}{y}})$. فقد ثبت أن ${\OLId{latin-looped}{R}}$ شاملة، خلافًا للّمة المتقدمة.
\end{proof}

ونطلق «الحساب الصادق» على النظرية
$\Setabs{!A}{\Sat{\Nat}{!A}}$، أي جميع جمل لغة الحساب التي تصدق
في التأويل القياسي.

\begin{cor}
لا يقبل الحساب الصادق القرار.
\end{cor}

%In Section~\olref{undefinability:truth:section} we will state a stronger
%result, due to Tarski.

\end{document}
